% GENERATED FILE. Edit canonical lessons, then run npm run generate:books.
% canonical-source-hash: fnv1a64:c0d44190c8b97c0f
% canonical-lessons: ZH-C463-danxiao, ZH-C463-kepa, ZH-W463-pa-fear, ZH-C463-xiaren, ZH-C463-fadou, ZH-C463-emeng, ZH-R463-review-458-463

\chapter{Scared}
\label{ch:zh-scared}
\hlchaptermodality{463}

\begin{chapteropening}
\textbf{By the end of this chapter:} \emph{I can say someone is timid, that something is frightening, ask someone not to scare me and talk about nightmares.}

Complete the last lesson of chapter 463: I can say someone is timid, that something is frightening, ask someone not to scare me and talk about nightmares.
\end{chapteropening}

\section[dǎnxiǎo]{dǎnxiǎo — timid; easily scared}

\label{lesson:ZH-C463-danxiao}



\begin{quote}
Before the new one: say the Mandarin for to give up hope, then the Mandarin for to fall out of love.
\end{quote}

\subsection*{You'll want to know: dǎnxiǎo}
\emph{dǎnxiǎo} — \textquotedblleft{}timid; easily scared\textquotedblright{}.

\begin{quote}
\emph{wǒ dìdi bǐ wǒ dǎnxiǎo} — My little brother scares more easily than I do.
\end{quote}

Its characters come later; for now it is a word for the ear.

Two third tones meet in this word: dǎn is said with a rising tone before xiǎo, and the spelling keeps both marks.

\begin{cousinweb}
\emph{dǎnxiǎo} is built from \emph{dǎn} 'the gall' and \zh{小} 'small', both inherited from Middle Chinese, \emph{tamX} and \emph{sjewX} (Guangyun, in Baxter's transcription), which glosses the first in the pair liver and gall. Courage was thought to sit in the gall, so a small gall is a timid heart; the dictionary does not say so itself.
\end{cousinweb}

\begin{grammarlens}[title={What you've built}]
One more word for saying how it was.
\end{grammarlens}

\subsection*{Your turn}
\begin{itemize}
\raggedright
  \item \emph{Say it:} \emph{dǎnxiǎo}
  \item \emph{Say it:} \emph{dǎnxiǎo}, once more
  \item \emph{Say it:} \emph{wǒ dìdi bǐ wǒ dǎnxiǎo} — My little brother scares more easily than I do.
  \item \emph{Recall:} say the Mandarin for to give up hope, then the Mandarin for to fall out of love, then say \emph{dǎnxiǎo} again
\end{itemize}

\begin{tcolorbox}[breakable,colback=teal!4,colframe=teal!35!black,title={Before you move on}]
What does dǎnxiǎo mean? (\textquotedblleft{}Timid; easily scared\textquotedblright{}.) Say it once more.
\end{tcolorbox}
\section[\texorpdfstring{\zh{可怕} (kěpà)}{可怕 (kěpà)}]{\zh{可怕} (kěpà) — frightening; awful}

\label{lesson:ZH-C463-kepa}



\begin{quote}
Before the new one: say the Mandarin for to fall out of love, then the Mandarin for easily scared.
\end{quote}

\subsection*{You'll want to know: \zh{可怕}}
\textbf{\zh{可怕}} — \emph{kěpà} — \textquotedblleft{}frightening; awful\textquotedblright{}.

\begin{quote}
\emph{zhège diànyǐng hěn kěpà} — This film is really frightening.
\end{quote}

The new character, in pieces: \zh{怕} — eight strokes in two pieces: on the left the heart piece, a dot, a second dot and a long vertical; on the right \zh{白} \emph{bái}, five strokes. The next lesson puts it in your hand.

\begin{cousinweb}
\textbf{\zh{可怕}} (\emph{kěpà}) is built from \zh{可} 'can' and \zh{怕} 'to fear', both inherited from Middle Chinese, \emph{khaX} and \emph{phaeH} (Guangyun, in Baxter's transcription), which glosses the first 'to permit' and the second 'to fear'. The dictionary lists \zh{怕} under a second reading, \emph{phaek}, glossed 'still and quiet', which is not this word.
\end{cousinweb}

\begin{grammarlens}[title={What you've built}]
One more word for saying how it was.
\end{grammarlens}

\subsection*{Your turn}
\begin{itemize}
\raggedright
  \item \emph{Say it:} \emph{kěpà}
  \item \emph{Say it:} \emph{kěpà}, once more
  \item \emph{Say it:} \emph{zhège diànyǐng hěn kěpà} — This film is really frightening.
  \item \emph{Recall:} say the Mandarin for to fall out of love, then the Mandarin for easily scared, then say \emph{kěpà} again
\end{itemize}

\begin{tcolorbox}[breakable,colback=teal!4,colframe=teal!35!black,title={Before you move on}]
What does \zh{可怕} mean? (\textquotedblleft{}Frightening; awful\textquotedblright{}.) Say it once more.
\end{tcolorbox}
\section[\texorpdfstring{\zh{怕} (pà)}{怕 (pà)}]{\zh{怕} — the heart piece, then \zh{白}}

\label{lesson:ZH-W463-pa-fear}



\begin{quote}
The word first, then the character inside it.

\begin{itemize}
\raggedright
  \item \emph{Recall:} say \emph{kěpà} — \textbf{R1}, one lesson back, before it decays
\end{itemize}

\begin{quote}
\zh{白} \zh{情}
\end{quote}

\emph{Read it:} each one aloud, with nothing said first — \emph{bái}, \emph{qíng}
\end{quote}

\subsection*{Script — the heart piece, then \zh{白}}
\hlblockfigure{\detokenize{figures/ZH-W463-pa-fear-filmstrip.pdf}}{How \zh{怕} is written, stroke by stroke}

\begin{quote}
\zh{怕}
\end{quote}

\textbf{Two pieces, left before right.} Count the pieces, not the strokes.

\textbf{The left piece, three strokes: the heart piece, as in \zh{情}.} A short dot on the left, a short dot on the right, then a long vertical down between them.

\textbf{The right piece, \zh{白}, five strokes.} A short stroke falling to the left at the top, a vertical down the left, a stroke that runs right and turns down, a short bar inside, then the bar across the foot.

3 + 5 = eight.

\subsection*{Your turn}
\begin{itemize}
\raggedright
  \item \emph{Write it:} \zh{怕}, all eight, with the model in front of you
  \item \emph{Cover:} write it again from memory, and count the pieces, not the strokes
\end{itemize}

\begin{tcolorbox}[breakable,colback=teal!4,colframe=teal!35!black,title={Before you move on}]
How many pieces does \zh{怕} have? (\textbf{Two: three strokes and five}.) Which written character is on the right? (\textbf{\zh{白}}.)

Source: \href{https://raw.githubusercontent.com/chanind/hanzi-writer-data/68d10a4b21150cae5e1ebbd223eed289cf32d90c/data/\%E6\%80\%95.json}{Hanzi Writer Data \zh{怕}, pinned stroke paths and medians}.
\end{tcolorbox}
\section[xiàrén]{xiàrén — scary; giving people a fright}

\label{lesson:ZH-C463-xiaren}



\begin{quote}
Before the new one: say the Mandarin for easily scared, then the Mandarin for frightening.
\end{quote}

\subsection*{You'll want to know: xiàrén}
\emph{xiàrén} — \textquotedblleft{}scary; giving people a fright\textquotedblright{}.

\begin{quote}
\emph{tiānsè hēi le wàibian yǒudiǎnr xiàrén} — It's getting dark, and it's a bit scary outside.
\end{quote}

Its characters come later; for now it is a word for the ear.

\begin{cousinweb}
\emph{xiàrén} is built from \emph{xià} and \zh{人} 'a person', both inherited from Middle Chinese, \emph{xaeH} and \emph{nyin} (Guangyun, in Baxter's transcription). The dictionary lists the first under two readings and glosses this one as the sound of laughing; the sense 'to frighten' is not given there.
\end{cousinweb}

\begin{grammarlens}[title={What you've built}]
One more word for saying how it was.
\end{grammarlens}

\subsection*{Your turn}
\begin{itemize}
\raggedright
  \item \emph{Say it:} \emph{xiàrén}
  \item \emph{Say it:} \emph{xiàrén}, once more
  \item \emph{Say it:} \emph{tiānsè hēi le wàibian yǒudiǎnr xiàrén} — It's getting dark, and it's a bit scary outside.
  \item \emph{Recall:} say the Mandarin for easily scared, then the Mandarin for frightening, then say \emph{xiàrén} again
\end{itemize}

\begin{tcolorbox}[breakable,colback=teal!4,colframe=teal!35!black,title={Before you move on}]
What does xiàrén mean? (\textquotedblleft{}Scary; giving people a fright\textquotedblright{}.) Say it once more.
\end{tcolorbox}
\section[fādǒu]{fādǒu — to tremble; to shiver}

\label{lesson:ZH-C463-fadou}



\begin{quote}
Before the new one: say the Mandarin for frightening, then the Mandarin for giving people a fright.
\end{quote}

\subsection*{You'll want to know: fādǒu}
\emph{fādǒu} — \textquotedblleft{}to tremble; to shiver\textquotedblright{}.

\begin{quote}
\emph{wàibian hěn lěng wǒ de shǒu zài fādǒu} — It's freezing outside, and my hands are shaking.
\end{quote}

Its characters come later; for now it is a word for the ear.

\begin{cousinweb}
\emph{fādǒu} is built from \zh{发} 'to set going' and \emph{dǒu} 'to shake', both inherited from Middle Chinese, \emph{pjot} and \emph{tuwX} (Guangyun, in Baxter's transcription), which glosses the first 'to start' and the second as shaking something out.
\end{cousinweb}

\begin{grammarlens}[title={What you've built}]
One more word for saying how it was.
\end{grammarlens}

\subsection*{Your turn}
\begin{itemize}
\raggedright
  \item \emph{Say it:} \emph{fādǒu}
  \item \emph{Say it:} \emph{fādǒu}, once more
  \item \emph{Say it:} \emph{wàibian hěn lěng wǒ de shǒu zài fādǒu} — It's freezing outside, and my hands are shaking.
  \item \emph{Recall:} say the Mandarin for frightening, then the Mandarin for giving people a fright, then say \emph{fādǒu} again
\end{itemize}

\begin{tcolorbox}[breakable,colback=teal!4,colframe=teal!35!black,title={Before you move on}]
What does fādǒu mean? (\textquotedblleft{}To tremble; to shiver\textquotedblright{}.) Say it once more.
\end{tcolorbox}
\section[èmèng]{èmèng — nightmare}

\label{lesson:ZH-C463-emeng}



\begin{quote}
Before the new one: say the Mandarin for giving people a fright, then the Mandarin for to tremble.
\end{quote}

\subsection*{You'll want to know: èmèng}
\emph{èmèng} — \textquotedblleft{}nightmare\textquotedblright{}.

\begin{quote}
\emph{wǒ zuò guo yí ge hěn kěpà de èmèng} — I once had a terrifying nightmare.
\end{quote}

Its characters come later; for now it is a word for the ear.

\begin{cousinweb}
\emph{èmèng} is built from \emph{è} and \emph{mèng} 'a dream', both inherited from Middle Chinese, \emph{ngak} and \emph{mjuwngH} (Guangyun, in Baxter's transcription). The dictionary gives the first only in the name of a year, and explains the second, after an older book, as the spirit wandering in sleep; the sense 'dreadful' for the first is not given there.
\end{cousinweb}

\begin{grammarlens}[title={What you've built}]
One more word for saying how it was.
\end{grammarlens}

\subsection*{Your turn}
\begin{itemize}
\raggedright
  \item \emph{Say it:} \emph{èmèng}
  \item \emph{Say it:} \emph{èmèng}, once more
  \item \emph{Say it:} \emph{wǒ zuò guo yí ge hěn kěpà de èmèng} — I once had a terrifying nightmare.
  \item \emph{Recall:} say the Mandarin for giving people a fright, then the Mandarin for to tremble, then say \emph{èmèng} again
\end{itemize}

\begin{tcolorbox}[breakable,colback=teal!4,colframe=teal!35!black,title={Before you move on}]
What does èmèng mean? (\textquotedblleft{}Nightmare\textquotedblright{}.) Say it once more.
\end{tcolorbox}
\section[(dialogue)]{Looking back: from How I Feel to Scared}

\label{lesson:ZH-R463-review-458-463}



\begin{quote}
Say the Mandarin for one's spirits and to feel for someone.
\end{quote}

\subsection*{Your turn}
\begin{itemize}
\raggedright
  \item \emph{Say it:} one's spirits, to feel for someone, thrilled, relaxed and at ease, fed up
  \item \emph{Say it:} lonely, a matter to deal with, affection between people, what is really on one's mind, hurt by unfair blame
  \item \emph{Say it:} down in the dumps, to sigh, embarrassed, to sympathize with, to envy
  \item \emph{Say it:} to hold in high regard, love and care for others, love between a couple, treasured, to look forward to
  \item \emph{Say it:} a pleasant surprise, beyond what anyone imagined, to hate deeply, to give up hope, to fall out of love
  \item \emph{Say it:} easily scared, frightening, giving people a fright, to tremble, a nightmare
  \item \emph{Read it:} \zh{情} \zh{爱} \zh{怕} — each one aloud, then a word that holds it
\end{itemize}

\begin{tcolorbox}[breakable,colback=teal!4,colframe=teal!35!black,title={Before you move on}]
One's spirits? (\emph{qíngxù}.) To feel for someone? (\textbf{\zh{心疼}}, \emph{xīnténg}.) Thrilled? (\emph{xīngfèn}.) Relaxed and at ease? (\emph{qīngsōng}.) Fed up? (\emph{fán}.) Lonely? (\emph{gūdān}.) A matter to deal with? (\textbf{\zh{事情}}, \emph{shìqing}.) Affection between people? (\textbf{\zh{感情}}, \emph{gǎnqíng}.) What is really on one's mind? (\textbf{\zh{心里话}}, \emph{xīnlihuà}.) Hurt by unfair blame? (\emph{wěiqu}.) Down in the dumps? (\emph{yùmèn}.) To sigh? (\emph{tànqì}.) Embarrassed? (\textbf{\zh{难为情}}, \emph{nánwéiqíng}.) To sympathize with? (\textbf{\zh{同情}}, \emph{tóngqíng}.) To envy? (\emph{xiànmù}.) To hold in high regard? (\emph{pèifú}.) Love and care for others? (\textbf{\zh{爱心}}, \emph{àixīn}.) Love between a couple? (\textbf{\zh{爱情}}, \emph{àiqíng}.) Treasured? (\textbf{\zh{心爱}}, \emph{xīn'ài}.) To look forward to? (\emph{pànwàng}.) A pleasant surprise? (\emph{jīngxǐ}.) Beyond what anyone imagined? (\textbf{\zh{意想不到}}, \emph{yìxiǎngbúdào}.) To hate deeply? (\emph{hèn}.) To give up hope? (\textbf{\zh{死心}}, \emph{sǐxīn}.) To fall out of love? (\textbf{\zh{变心}}, \emph{biànxīn}.) Easily scared? (\emph{dǎnxiǎo}.) Frightening? (\textbf{\zh{可怕}}, \emph{kěpà}.) Giving people a fright? (\emph{xiàrén}.) To tremble? (\emph{fādǒu}.) A nightmare? (\emph{èmèng}.)
\end{tcolorbox}
